% ============================================================================
% DRAFTING HEADER (strip before deposit)
% Drafted by Claude Fable 5.1, October 23, 2026, at the author's direction ("please go ahead and write the
% psychology paper now"), as the sixth extraction of the corpus and the first from the psychological material of the
% long Aethic paper: the section "The Aethic Mindset and Autistic Savants" and its subsections, the "it toggles"
% paragraph of "Aethic Reasoning Against the Modal Structure of Human Language", and the author's notes on his own
% record (the source file 1_Aethic_Reasoning.tex, md5 41fb70e5c55830eb139c50dc179f3009, line numbers below refer to it).
% The plan is branch_symmetry_psych_scope.md, third version (Oct 22, 2026), revised through two referee reports; the
% paper itself was revised to a third report (R848, Oct 23, 2026). The author's passages are quoted verbatim and marked
% (the two principles' "So, Branch Principle:" and "So, Principle of Symmetry:" sentences are in the Oct 22 source
% version only, not in the Oct 12 repository copy: the deposited long paper must carry the Oct 22 text); the formal
% apparatus, the design and the simulation are the extraction's. The simulation is norm_cost_sim.py, version 1.2,
% whose output accompanies the paper.
% Author-side card: the references marked VERIFY were written from memory of the works and need a check against the
% volumes before deposit; strip these working comments and this header before deposit.
% ============================================================================
\begin{filecontents}[overwrite]{norm_cost_refs.bib}
@misc{Aethus, author={Benander, Ajax}, title={Aethic Reasoning: A Comprehensive Solution to the Quantum Measurement Problem}, year={2024}, month = {11}, url={https://philpapers.org/rec/BENARA-5}, note={Manuscript; the complete development of the framework, successive deposits November 2024--2026; also at \url{https://philsci-archive.pitt.edu/id/eprint/25713}}}
@article{Milton2012,
  author = {Milton, Damian E. M.},
  title = {On the ontological status of autism: the `double empathy problem'},
  journal = {Disability \& Society}, volume = {27}, number = {6}, pages = {883--887}, year = {2012},
  doi = {10.1080/09687599.2012.710008}}
@article{Milton2022,
  author = {Milton, Damian and Gurbuz, Emine and L{\'o}pez, Beatriz},
  title = {The `double empathy problem': Ten years on},
  journal = {Autism}, volume = {26}, number = {8}, pages = {1901--1903}, year = {2022},
  doi = {10.1177/13623613221129123}}
% VERIFY at merge: fields from memory.
@article{Pellicano2022,
  author = {Pellicano, Elizabeth and den Houting, Jacquiline},
  title = {Annual Research Review: Shifting from `normal science' to neurodiversity in autism science},
  journal = {Journal of Child Psychology and Psychiatry}, volume = {63}, number = {4}, pages = {381--396}, year = {2022},
  doi = {10.1111/jcpp.13534}}
@article{DavisCrompton2021,
  author = {Davis, Rachael and Crompton, Catherine J.},
  title = {What do new findings about social interaction in autistic adults mean for neurodevelopmental research?},
  journal = {Perspectives on Psychological Science}, volume = {16}, number = {3}, pages = {649--653}, year = {2021},
  doi = {10.1177/1745691620958010}}
@article{Crompton2020,
  author = {Crompton, Catherine J. and Ropar, Danielle and Evans-Williams, Claire V. M. and Flynn, Emma G. and Fletcher-Watson, Sue},
  title = {Autistic peer-to-peer information transfer is highly effective},
  journal = {Autism}, volume = {24}, number = {7}, pages = {1704--1712}, year = {2020},
  doi = {10.1177/1362361320919286}}
@article{Crompton2025,
  author = {Crompton, Catherine J. and Foster, Sarah J. and Wilks, Charlotte E. H. and Dodd, Michelle and Efthimiou, Themis N. and Ropar, Danielle and Sasson, Noah J. and Lages, Martin and Fletcher-Watson, Sue},
  title = {Information transfer within and between autistic and non-autistic people},
  journal = {Nature Human Behaviour}, volume = {9}, number = {7}, pages = {1488--1500}, year = {2025},
  doi = {10.1038/s41562-025-02163-z}}
@article{Foster2025,
  author = {Foster, Sarah J. and Ackerman, Robert A. and Wilks, Charlotte E. H. and Dodd, Michelle and Calderon, Rachel and Ropar, Danielle and Fletcher-Watson, Sue and Crompton, Catherine J. and Sasson, Noah J.},
  title = {Rapport in same and mixed neurotype groups of autistic and non-autistic adults},
  journal = {Autism}, volume = {29}, number = {7}, year = {2025},
  doi = {10.1177/13623613251320444}}
@article{Sasson2017,
  author = {Sasson, Noah J. and Faso, Daniel J. and Nugent, Jack and Lovell, Sarah and Kennedy, Daniel P. and Grossman, Ruth B.},
  title = {Neurotypical peers are less willing to interact with those with autism based on thin slice judgments},
  journal = {Scientific Reports}, volume = {7}, pages = {40700}, year = {2017},
  doi = {10.1038/srep40700}}
@misc{Calderoli2026,
  author = {Calderoli, Enrique and Figueiredo, Tiago and Varriale, Maria Cristina and Kapczinski, Fl{\'a}vio},
  title = {A distinct communication strategies model of the double empathy problem},
  note = {arXiv:2602.02562, version 2, revised October 1, 2026}, year = {2026}}
@article{Cronbach1957,
  author = {Cronbach, Lee J.},
  title = {The two disciplines of scientific psychology},
  journal = {American Psychologist}, volume = {12}, number = {11}, pages = {671--684}, year = {1957},
  doi = {10.1037/h0043943}}
@book{CronbachSnow1977,
  author = {Cronbach, Lee J. and Snow, Richard E.},
  title = {Aptitudes and Instructional Methods: A Handbook for Research on Interactions},
  publisher = {Irvington}, address = {New York}, year = {1977}}
@book{Kenny2006,
  author = {Kenny, David A. and Kashy, Deborah A. and Cook, William L.},
  title = {Dyadic Data Analysis},
  publisher = {Guilford Press}, address = {New York}, year = {2006}}
@incollection{KennyLaVoie1984,
  author = {Kenny, David A. and La Voie, Lawrence},
  title = {The social relations model},
  booktitle = {Advances in Experimental Social Psychology}, editor = {Berkowitz, Leonard}, volume = {18}, pages = {141--182},
  publisher = {Academic Press}, address = {Orlando}, year = {1984}}
@article{ImaiKeeleTingley2010,
  author = {Imai, Kosuke and Keele, Luke and Tingley, Dustin},
  title = {A general approach to causal mediation analysis},
  journal = {Psychological Methods}, volume = {15}, number = {4}, pages = {309--334}, year = {2010},
  doi = {10.1037/a0020761}}
@article{ImaiKeeleYamamoto2010,
  author = {Imai, Kosuke and Keele, Luke and Yamamoto, Teppei},
  title = {Identification, inference and sensitivity analysis for causal mediation effects},
  journal = {Statistical Science}, volume = {25}, number = {1}, pages = {51--71}, year = {2010},
  doi = {10.1214/10-STS321}}
@article{Lakens2017,
  author = {Lakens, Dani{\"e}l},
  title = {Equivalence tests: A practical primer for $t$ tests, correlations, and meta-analyses},
  journal = {Social Psychological and Personality Science}, volume = {8}, number = {4}, pages = {355--362}, year = {2017},
  doi = {10.1177/1948550617697177}}
@article{Lakens2018,
  author = {Lakens, Dani{\"e}l and Scheel, Anne M. and Isager, Peder M.},
  title = {Equivalence testing for psychological research: A tutorial},
  journal = {Advances in Methods and Practices in Psychological Science}, volume = {1}, number = {2}, pages = {259--269}, year = {2018},
  doi = {10.1177/2515245918770963}}
@book{Senn2002,
  author = {Senn, Stephen},
  title = {Cross-over Trials in Clinical Research},
  edition = {2}, publisher = {Wiley}, address = {Chichester}, year = {2002}}
@article{Stone1977,
  author = {Stone, Mervyn},
  title = {An asymptotic equivalence of choice of model by cross-validation and {A}kaike's criterion},
  journal = {Journal of the Royal Statistical Society, Series B}, volume = {39}, number = {1}, pages = {44--47}, year = {1977}}
@article{Shao1993,
  author = {Shao, Jun},
  title = {Linear model selection by cross-validation},
  journal = {Journal of the American Statistical Association}, volume = {88}, number = {422}, pages = {486--494}, year = {1993},
  doi = {10.1080/01621459.1993.10476299}}
@incollection{HartStaveland1988,
  author = {Hart, Sandra G. and Staveland, Lowell E.},
  title = {Development of {NASA-TLX} ({Task Load Index}): Results of empirical and theoretical research},
  booktitle = {Human Mental Workload}, editor = {Hancock, Peter A. and Meshkati, Najmedin},
  publisher = {North-Holland}, address = {Amsterdam}, pages = {139--183}, year = {1988}}
@article{Nosek2018,
  author = {Nosek, Brian A. and Ebersole, Charles R. and DeHaven, Alexander C. and Mellor, David T.},
  title = {The preregistration revolution},
  journal = {Proceedings of the National Academy of Sciences}, volume = {115}, number = {11}, pages = {2600--2606}, year = {2018},
  doi = {10.1073/pnas.1708274114}}
@book{APA2013,
  author = {{American Psychiatric Association}},
  title = {Diagnostic and Statistical Manual of Mental Disorders},
  edition = {5}, publisher = {American Psychiatric Publishing}, address = {Arlington, VA}, year = {2013}}
@article{SattathTversky1977,
  author = {Sattath, Shmuel and Tversky, Amos},
  title = {Additive similarity trees},
  journal = {Psychometrika}, volume = {42}, number = {3}, pages = {319--345}, year = {1977},
  doi = {10.1007/BF02293654}}
@article{Tversky1977,
  author = {Tversky, Amos},
  title = {Features of similarity},
  journal = {Psychological Review}, volume = {84}, number = {4}, pages = {327--352}, year = {1977},
  doi = {10.1037/0033-295X.84.4.327}}
@article{Litman2025,
  author = {Litman, Aviya and Sauerwald, Natalie and Green Snyder, LeeAnne and Foss-Feig, Jennifer and Park, Christopher Y. and Hao, Yun and Dinstein, Ilan and Theesfeld, Chandra L. and Troyanskaya, Olga G.},
  title = {Decomposition of phenotypic heterogeneity in autism reveals underlying genetic programs},
  journal = {Nature Genetics}, volume = {57}, number = {7}, pages = {1611--1619}, year = {2025},
  doi = {10.1038/s41588-025-02224-z}}
@article{WilliamsGotham2021,
  author = {Williams, Zachary J. and Gotham, Katherine O.},
  title = {Improving the measurement of alexithymia in autistic adults: a psychometric investigation of the 20-item {Toronto Alexithymia Scale} and generation of a general alexithymia factor score using item response theory},
  journal = {Molecular Autism}, volume = {12}, pages = {56}, year = {2021},
  doi = {10.1186/s13229-021-00463-5}}
% VERIFY at merge (the three below): fields from memory; the long paper's refs.bib carries the canonical entries.
@article{Hull2017,
  author = {Hull, Laura and Petrides, K. V. and Allison, Carrie and Smith, Paula and Baron-Cohen, Simon and Lai, Meng-Chuan and Mandy, William},
  title = {`Putting on my best normal': Social camouflaging in adults with autism spectrum conditions},
  journal = {Journal of Autism and Developmental Disorders}, volume = {47}, number = {8}, pages = {2519--2534}, year = {2017},
  doi = {10.1007/s10803-017-3166-5}}
@article{Hull2019,
  author = {Hull, Laura and Mandy, William and Lai, Meng-Chuan and Baron-Cohen, Simon and Allison, Carrie and Smith, Paula and Petrides, K. V.},
  title = {Development and validation of the {Camouflaging Autistic Traits Questionnaire} ({CAT-Q})},
  journal = {Journal of Autism and Developmental Disorders}, volume = {49}, number = {3}, pages = {819--833}, year = {2019},
  doi = {10.1007/s10803-018-3792-6}}
@article{Pearson2021,
  author = {Pearson, Amy and Rose, Kieran},
  title = {A conceptual analysis of autistic masking: Understanding the narrative of stigma and the illusion of choice},
  journal = {Autism in Adulthood}, volume = {3}, number = {1}, pages = {52--60}, year = {2021},
  doi = {10.1089/aut.2020.0043}}
@article{Hajek2007,
  author = {H{\'a}jek, Alan},
  title = {The reference class problem is your problem too},
  journal = {Synthese}, volume = {156}, number = {3}, pages = {563--585}, year = {2007},
  doi = {10.1007/s11229-006-9138-5}}
@article{Robinson1950,
  author = {Robinson, W. S.},
  title = {Ecological correlations and the behavior of individuals},
  journal = {American Sociological Review}, volume = {15}, number = {3}, pages = {351--357}, year = {1950},
  doi = {10.2307/2087176}}
@article{BottemaBeutel2021,
  author = {Bottema-Beutel, Kristen and Kapp, Steven K. and Lester, Jessica Nina and Sasson, Noah J. and Hand, Brittany N.},
  title = {Avoiding ableist language: Suggestions for autism researchers},
  journal = {Autism in Adulthood}, volume = {3}, number = {1}, pages = {18--29}, year = {2021},
  doi = {10.1089/aut.2020.0014}}
\end{filecontents}
\documentclass[11pt]{article}
\usepackage[margin=1.05in]{geometry}
\usepackage{amsmath, amssymb, amsthm}
\usepackage[hidelinks]{hyperref}
\usepackage{microtype}
\newcommand{\principleautorefname}{Principle}
\newcommand{\propositionautorefname}{Proposition}
\newcommand{\definitionautorefname}{Definition}
\newcommand{\conjectureautorefname}{Conjecture}
\newcommand{\hypothesisautorefname}{Hypothesis}
\newcommand{\corollaryautorefname}{Corollary}
\renewcommand{\sectionautorefname}{Section}
\renewcommand{\subsectionautorefname}{Section}
\usepackage{graphicx}
\usepackage{booktabs}
\usepackage{array}
\usepackage{enumitem}
\usepackage[style=numeric, sorting=none, backend=bibtex]{biblatex}
\addbibresource{norm_cost_refs.bib}

\newtheorem{theorem}{Theorem}
\newtheorem{proposition}{Proposition}
\newtheorem{principle}{Principle}
\newtheorem{corollary}{Corollary}
\theoremstyle{definition}
\newtheorem{definition}{Definition}
\newtheorem{conjecture}{Conjecture}
\newtheorem{hypothesis}{Hypothesis}

\newcommand{\pos}[1]{\left[#1\right]_{+}}

\title{Reference, Norms and the Cost of Interaction:\\ A Norm-Setting Model of Autistic--Non-Autistic Communication, with a Discriminating Design}
\author{Ajax Benander\thanks{Email: \href{mailto:akb9408@rit.edu}{akb9408@rit.edu}, Website: \url{https://aethus.org}. This paper develops, as a self-contained proposal in the psychology of communication, two principles stated in a larger framework \cite{Aethus}: the Branch Principle and the Principle of Symmetry, with the observation that a difference which switches with context refutes an unconditional claim of incapacity. The framework's metaphysics plays no role in the argument; what is carried over is one piece of bookkeeping, the treatment of a diagnostic label as a partial information state, explained where it enters.}}
\date{October 23, 2026 --- working draft}

\begin{document}
\maketitle

\begin{abstract}
How much of the cost of an interaction between an autistic and a non-autistic person depends on the norms that govern it, and when can changing those norms change, or reverse, which partner pays? The paper gives the question a formal setting. A diagnostic label is treated as a partial information state, which fixes what a claim about ``the autistic'' may and may not say; a profile over competences makes deficit verdicts reference-relative, and reversibility, at a stated measurement margin, is proved to hold exactly under a condition of mutual advantage, with ``balance'' kept as a conjecture. A cost model is then stated: effort is a participant baseline plus a symmetric mismatch between the participant's communicative style and the style the governing norms presuppose, plus a one-sided shortfall against the competence the norms require. Its predictions are derived exactly: the group gap under given norms, its sensitivity to a change of norms, from which the baseline cancels within persons, and the condition under which the gap reverses, namely that the norm-induced reduction exceed the whole initial gap, not the baseline alone. The model is one member of a family that includes a general protocol benefit and partner mismatch as restrictions; the double-empathy tradition's formal model is not a restriction of it, and a norm-dependent extension of that model can be separated from the norm model only under mediation assumptions the design cannot enforce, so that analysis is stated as assumption-dependent and the mediator's manipulation is left to a later study. A within-dyad crossover is designed in which the same mixed pairs work under two explicit rule sets, matched on specified incidental features and differing in the profile they presuppose, with profiles measured before any interaction. A simulation at the planned size recovers the parameters, selects the generating model in most replications by the Bayesian information criterion on an exact likelihood, predicts the relative individual effects of a withheld third rule set for held-out dyads to the noise floor, exhibits the reversal arithmetic, and reports which of these conclusions survive measurement error, correlated baselines, correlated predictors, unequal variances and smaller effects. What existing results constrain is stated from the primary sources. Branch geometry and linguistic masking are extensions, not premises.
\end{abstract}

\section{Introduction}
\label{sec:intro}

\subsection{The question}

When an autistic and a non-autistic person interact and the interaction goes badly, the received description locates the difficulty in one of them. The double empathy problem \cite{Milton2012} relocated it to the pair: a mutual failure of understanding between people whose ways of experiencing and communicating differ, not a one-sided deficit, a relocation with a decade of empirical work behind it \cite{Milton2022, DavisCrompton2021, Pellicano2022}. This paper takes one further step and asks about the third party present in every such interaction: the norms that govern it. Any interaction runs under rules about how a request is made, how a turn is yielded, how much is said outright and how much left to be inferred, and those rules presuppose a communicator. The question is how much of the interaction's cost is set by that presupposition, and when changing the rules changes, or reverses, which partner pays it.

The question has three parts, and the paper keeps them apart because they have different standing.
\begin{itemize}[leftmargin=*]
\item \emph{A definitional part.} A verdict of deficit is relative to a reference. This is a matter of definition, and it is proved below in the only sense a definition can be proved, by exhibiting the quantifier the verdict hides.
\item \emph{A formal part.} Under a cost model with a stated structure, the group gap in effort under given norms, its change under new norms and the condition for its reversal follow exactly, and the reversal condition is stricter than it is usually taken to be.
\item \emph{An empirical part.} Whether effort tracks the mismatch between a person and the norms in the way the model says is a hypothesis, and the paper's main work is a design that can tell it from its strongest rivals.
\end{itemize}

\subsection{What the paper adds}

The contribution is operational. Interactions between differently situated people have an established theoretical setting in the study of aptitude-by-treatment interactions \cite{Cronbach1957, CronbachSnow1977}: who benefits from a condition depends jointly on the person and the condition, and context-dependent advantages and their reversals are at home there. The double-empathy literature has a formal model of its own, a feedback model of channel preference in which each partner's registered empathy depends on the partner's verbal and nonverbal output weighted by the first partner's preference \cite{Calderoli2026}. What this paper adds is narrower than either and not contained in them: a prediction of how an externally imposed communication norm distributes the cost of adaptation between the two partners of a dyad, from profiles measured before the interaction and norms coded independently of its outcome, beyond the effect of who the partner is. The model that carries the prediction is one of a family stated together in \autoref{sec:family}, each rival a parameter restriction, and no tradition is identified with a restriction.

\subsection{Where the model comes from}
\label{sec:source}

The paper develops two principles the author stated in a larger work \cite{Aethus}, in a section on the relation between cognitive style and the conventions of a dominant group; Appendix~\ref{app:source} gives the passages in full, with their provenance. The \emph{Branch Principle} holds that ``ontologically neurodivergent people are better understood as a collection of categorical branches and `allistic' is a single branch, such that two random branches are typically as different from one another as they each are to a third''. The \emph{Principle of Symmetry} holds that ``you can always redraw the conventional category boundaries in psychology so that allistic and a given neurodivergent branch `trade places phenomenologically'{}''. Of what the passages say, the paper separates four things. Reference-relativity is definitional (\autoref{sec:symmetry}). Reversibility, the ``trade places'' claim, is an exact consequence of a premise the passage states, that ``each branch comes with a different set of cognitive intuitions'', once that premise is read as mutual advantage at the measurement's margin; without it, it fails, and the failure case is exhibited. ``Roughly half-and-half'' is a conjecture, and nothing below rests on it. And the remark that autistic people ``are described as taking a script to allistic conversations'', when ``if the shoe were on the other foot'' allistic people would be, is the seed of the cost model: the labour an interaction extracts depends on whose conventions it runs on, so that a change of conventions moves the labour. That is the hypothesis the paper is built to test. The Branch Principle itself, a claim about the geometry of cognitive profiles, is placed outside this paper as an extension with its own test (\autoref{sec:extensions}).

Two commitments govern the writing. First, the paper's claims about ``the autistic'' obey a scope rule stated in \autoref{sec:labels}, which distinguishes what a label entails of every bearer, what is a prediction about a distribution, and what is attributed to an individual. Second, the language follows the recommendations made to autism researchers on ableist language \cite{BottemaBeutel2021}: ``autistic'' and ``non-autistic'' are used as the groups' names, ``allistic'' only inside the source's quotations, and no competence difference is called a deficit except relative to a stated reference.

\subsection{Plan}

Sections~\ref{sec:labels} to \ref{sec:toggle} are compact preliminaries: a label as a partial information state and the scope rule; the Principle of Symmetry made exact; what a difference that switches with context does and does not show. Section~\ref{sec:model} gives the cost model and derives its predictions, Section~\ref{sec:family} the model family and what can and cannot separate its members, Section~\ref{sec:design} the study, Section~\ref{sec:simulation} the simulation with its robustness checks, Section~\ref{sec:constraints} what the published results already constrain, Section~\ref{sec:extensions} what lies outside the paper, and Section~\ref{sec:discussion} the discussion. Appendix~\ref{app:estimators} states the simulation's likelihood, Appendix~\ref{app:prereg} the preregistration items, and Appendix~\ref{app:source} the source passages.

\section{A label as a partial information state}
\label{sec:labels}

A diagnostic label carries what it records and no more. The larger framework calls such a state an Aethus and any fuller description of the same person a proper child of it \cite{Aethus}; this paper says \emph{partial information state} and \emph{refinement}. The point of the vocabulary is not metaphysical. It is that a label is an object with a definite content, and that a claim about everyone who bears it is licensed by that content or not at all.

Three kinds of claim about the bearers of a label are therefore distinguished.
\begin{enumerate}[label=(\roman*), leftmargin=*]
\item \emph{Definitional entailment.} What the label records holds of every bearer. The diagnostic criteria for autism spectrum disorder in the fifth edition of the \emph{Diagnostic and Statistical Manual} require all three of the social-communication features and at least two of the four restricted and repetitive features, currently or by history \cite{APA2013}. So ``every bearer shows at least two of the four'' is entailed. ``Every bearer has sensory differences'' is not: sensory reactivity is one of the four, and the label does not record which two.
\item \emph{Group-level statistical prediction.} A claim about the distribution of bearers. It is true or false of the group, and attributing it to each member is the ecological fallacy \cite{Robinson1950}.
\item \emph{Individual attribution.} A claim about a particular bearer, which needs information about that bearer.
\end{enumerate}

\begin{definition}[Scope rule]
\label{def:scope}
A claim about every bearer of a label is licensed in two ways only: the label's definition entails it, or independent evidence establishes it for every bearer. A claim that needs information the label does not carry is made of a refinement that carries it, or not made.
\end{definition}

\noindent The rule is the discipline the rest of the paper is written under. It does not forbid universal claims; it says where they come from.

\paragraph{The source's illustration.} The section the principles come from gives one worked instance of a claim that violates the rule (Appendix~\ref{app:source}): for any trait some autistic branch can be found on which it presents as infantile, which seems to confirm that ``autistic people are infantile in this way'', and the same inference run on a random partition of the branches, the allistic one among them, would convict allistic people of anything. With $I(b, \tau)$ reading, as the source's formalization has it, ``$\tau$ presents as infantile on $b$ relative to a chosen reference branch'', observation supplies $\forall \tau\, \exists b\; I(b, \tau)$ and the category-level conclusion needs $\exists b\, \forall \tau\; I(b, \tau)$: a quantifier shift, with the witness fixed per demonstration and silently promoted, whose confirmations are a function of the partition chosen, the reference class problem \cite{Hajek2007} in categorical form. The logic is elementary; what the illustration shows is that a specific and common inference makes the move.

\section{The Principle of Symmetry made exact}
\label{sec:symmetry}

\subsection{Profiles, references and deficit sets}

Let $T$ be a finite inventory of competences and let each branch $b$ have a profile $x_b \colon T \to \mathbb{R}$, a score on each competence on a common scale. The scale carries a measurement margin $\delta \ge 0$, the smallest difference the instrument resolves or the analyst is prepared to call a difference.

\begin{definition}[Deficit set]
\label{def:deficit}
The deficit set of branch $b$ relative to a reference branch $r$ at margin $\delta$ is
\[ D_r(b;\delta) \;=\; \{\, \tau \in T : x_b(\tau) < x_r(\tau) - \delta \,\}. \]
An inventory $S \subseteq T$ shows $b$ in deficit relative to $r$ when $S \subseteq D_r(b;\delta)$ and $S \neq \emptyset$.
\end{definition}

\noindent Three statements follow, of three different kinds.

\paragraph{(S1) Reference-relativity, a definition.} A deficit verdict is a three-place relation among a branch, a reference and a competence, and $D_r(b;\delta)$ has no reference-free extension: there is no set ``the deficits of $b$'' until $r$ is fixed. Choosing another reference makes another comparison, not a correction of the first. This is the content of the source's formalization, ``the reference-relativity being the crux rather than a modeling convenience''.

\paragraph{(S2) Reversibility at margin $\delta$, an exact consequence.}
\begin{proposition}[Reversibility]
\label{prop:reversibility}
For branches $b$ and $r$ and margin $\delta \ge 0$, the following are equivalent:
\begin{enumerate}[label=(\alph*), leftmargin=*, nosep]
\item some inventory shows $b$ in deficit relative to $r$, and some inventory shows $r$ in deficit relative to $b$;
\item there are competences $\tau$ and $\sigma$ with $x_b(\tau) - x_r(\tau) > \delta$ and $x_r(\sigma) - x_b(\sigma) > \delta$.
\end{enumerate}
Call (b) \emph{$\delta$-mutual advantage}.
\end{proposition}
\begin{proof}
(a) says that $D_r(b;\delta)$ and $D_b(r;\delta)$ are both non-empty. $\sigma \in D_r(b;\delta)$ is $x_r(\sigma) - x_b(\sigma) > \delta$, and $\tau \in D_b(r;\delta)$ is $x_b(\tau) - x_r(\tau) > \delta$, which is (b).
\end{proof}

\noindent Under $\delta$-mutual advantage, then, the ``trade places'' of the Principle of Symmetry (\autoref{sec:source}) holds in this sense: restricting the inventory to $D_r(b;\delta)$ shows $b$ in deficit everywhere, and restricting it to $D_b(r;\delta)$ shows $r$ in deficit everywhere, so that the verdict is reversed by exchanging the reference and the inventory. It holds exactly when the premise holds and fails exactly when it does not. Two remarks fix the premise.
\begin{itemize}[leftmargin=*]
\item \emph{Non-domination is not enough above $\delta = 0$.} At $\delta = 0$, (b) is the statement that neither profile is at or above the other on every competence. At a positive margin it is stronger. Take $T = \{\tau_1, \tau_2\}$, $x_b = (1.05, 1)$, $x_r = (1, 1.05)$ and $\delta = 0.1$: neither profile dominates, yet both deficit sets are empty, and no inventory shows either in deficit.
\item \emph{What is testable.} Whether a non-reference branch exceeds the reference somewhere by more than the margin is a question about measured profiles, and it is answered competence by competence. The source's own case, two competences at opposite developmental ages (Appendix~\ref{app:source}), illustrates the premise's shape; it is offered as motivation and not as a measurement on commensurate scales, and it carries no evidential weight here.
\end{itemize}

\paragraph{(S3) Balance, a conjecture.}
\begin{conjecture}[Balance]
\label{conj:balance}
For the reference branch and a given non-reference branch, the deficit sets $D_r(b;\delta)$ and $D_b(r;\delta)$ carry comparable weight under a natural weighting of the competences.
\end{conjecture}
\noindent This is the source's ``roughly half-and-half''. It needs what \autoref{prop:reversibility} does not: a weighting of competences under which the two sets can be compared, and exchangeability of the profile differences under that weighting. No weighting is natural until the inventory's granularity is settled, since splitting one competence into five multiplies its weight, and the paper draws nothing from the conjecture. Reversibility in existence, (S2), is what the cost model below uses, and it uses even that only as motivation.

\subsection{Why mutual advantage is not a premise of what follows}

The cost model of \autoref{sec:model} predicts that a change of norms moves effort, and it does so whether or not the two groups stand in $\delta$-mutual advantage. A participant can gain from a change of norms while scoring lower on every measured competence, if the new norms presuppose less of what that participant lacks; and a participant who exceeds the reference on some competence gains nothing from norms that do not draw on it. The principles of \autoref{sec:source} motivate the comparison the model makes; they are not its premises, and the model's test does not test them. That separation is deliberate, since it lets the empirical work proceed without first settling the geometry of cognitive profiles, which is the Branch Principle's territory (\autoref{sec:extensions}).

\section{Context sensitivity: what toggling shows}
\label{sec:toggle}

The source makes one further observation, in a different section and about a different convention, which bears on how any difference between groups should be read: ``A limitation does not toggle.'' An adult who switches between two modes of a convention within a day, by context, shows that ``the convention is a mode the platform runs, selectable, not a ceiling the platform imposes'' (the passage is in Appendix~\ref{app:source}). The observation is stated for one person switching between two modes by context, and against one rival, a fixed limitation. Against that rival it is decisive, and the paper takes exactly that much from it.

\begin{proposition}[What a toggle refutes]
\label{prop:toggle}
Let $P$ be a performance and $C$ a set of conditions. If $P$ is observed under some $c \in C$, the claim that the person cannot produce $P$ under any condition in $C$ is false.
\end{proposition}
\noindent The proposition is immediate, and its value is in what it does not say. Two converses that would make the toggle a general test of capacity against convention do not hold.
\begin{itemize}[leftmargin=*]
\item \emph{Persistence does not identify capacity.} A convention can persist across every instrument that shares its assumptions, and most instruments in a field do.
\item \emph{Change does not identify convention.} A stable capacity yields different performance under different demands: with more time, clearer instructions, fewer simultaneous demands, another format. The performance changes while the constraint remains consequential.
\end{itemize}
The toggle is therefore a test of \emph{context sensitivity}: does the difference change across conditions? A change refutes the unconditional claim of incapacity and opens a question of mechanism, which further experiments answer by varying demand, time, load and format separately from the governing norm. In the study of \autoref{sec:design} the two conditions differ in the governing norm, are matched on specified incidental features, and have their remaining differences in training and practice measured, which is what lets a change be read as the norm's up to those measured differences.

Applying the observation to differences between groups is this paper's extension, and it is labelled as such. The participant-specific baseline $a_i$ of the model below is accordingly a baseline over the conditions studied and not a measure of capacity, and the components of it that might vary with the conditions, sensory load and fatigue, are held fixed or measured in the design.

\section{The cost of interaction under governing norms}
\label{sec:model}

\subsection{Profiles and norms}

Each participant $i$ has two profiles, both measured before any interaction.
\begin{itemize}[leftmargin=*]
\item A \emph{style profile} $s_i \in \mathbb{R}^{p}$: the participant's communicative preferences on dimensions such as directness, literalness, turn-taking and channel, scored on common scales.
\item A \emph{competence profile} $k_i \colon T \to \mathbb{R}$ over the abilities the interaction draws on, the profile of \autoref{sec:symmetry} measured for individuals.
\end{itemize}
A governing norm $n$, a rule set under which the interaction runs, fixes two things.
\begin{itemize}[leftmargin=*]
\item A \emph{presupposed style} $s^{*}_n \in \mathbb{R}^{p}$: the communicator the rules are written for, coded from the rules' text by raters who do not see the outcomes.
\item A \emph{competence requirement} $\theta_n \colon T \to \mathbb{R}$ with weights $w_n \colon T \to [0, 1]$: how much of each competence the rules call on, and how heavily.
\end{itemize}

\subsection{The cost model}

\begin{hypothesis}[The cost model]
\label{hyp:cost}
Participant $i$'s effort under norm $n$ is
\begin{equation}
\label{eq:cost}
c_i(n) \;=\; a_i \;+\; m_i(n) \;+\; \varepsilon_i(n), \qquad
m_i(n) \;=\; \lambda_s\, d\big(s_i, s^{*}_n\big) \;+\; \lambda_k \sum_{\tau \in T} w_n(\tau)\, \pos{\theta_n(\tau) - k_i(\tau)},
\end{equation}
with $\pos{u} = \max(0, u)$, $d$ a distance on style profiles, $\lambda_s, \lambda_k \ge 0$, $a_i$ a participant-specific baseline over the conditions studied and $\varepsilon_i(n)$ zero-mean noise.
\end{hypothesis}

\noindent The two terms of the mismatch $m_i(n)$ have different shapes, and the difference is the point.
\begin{itemize}[leftmargin=*]
\item \emph{The style term is symmetric.} Departing from the presupposed style in either direction costs effort: a direct communicator under indirect rules and an indirect one under direct rules both adapt. A distance is the right form for a mismatch of styles, where neither pole is better.
\item \emph{The competence term is one-sided.} Falling short of a requirement costs effort; exceeding it costs nothing. A participant whose competence exceeds what the rules call on is not penalized for the excess, which a symmetric distance on competences would wrongly do.
\item \emph{The baseline is a baseline.} $a_i$ collects whatever is constant across the conditions studied: capacity, in part, but also sensory load, fatigue and anything else the conditions share. It is not identified as capacity, and the model does not call it that.
\end{itemize}

\paragraph{Status.} Equation~\eqref{eq:cost} is a hypothesis about effort, which the paper tests. Reference-relativity, (S1), makes the comparison between a participant and a norm well defined; the hypothesis says how effort tracks it. Neither the Branch Principle nor reference-relativity entails the equation, and nothing in the framework it comes from does either.

\paragraph{Outcomes.} Effort is the primary outcome, measured per participant and per condition by a subjective workload scale, the mental-demand and effort items of the NASA Task Load Index \cite{HartStaveland1988} adapted to conversation, and by one behavioural index fixed in advance. Task performance and rapport are secondary outcomes, each with a model of its own; the model does not derive rapport from costs, since ``depends on'' yields no prediction. Expected and perceived partner empathy are measured in each condition, for the mediation analysis of \autoref{sec:family}.

\subsection{The predictions, exactly}
\label{sec:predictions}

Write $G(n)$ for the expected group gap in effort under norm $n$, the minority group's mean less the majority group's,
\begin{equation}
\label{eq:gap}
G(n) \;=\; \Delta a \;+\; \Delta M(n),
\end{equation}
with $\Delta a$ the baseline gap and $\Delta M(n)$ the gap in mean mismatch, both signed in the same direction.

\paragraph{P1, the gap under the default norms $n_0$.} $G(n_0) = \Delta a + \Delta M(n_0)$. Under norms that presuppose the majority's style and draw on its competences, $\Delta M(n_0) > 0$ whenever the minority's profiles lie farther from $s^{*}_{n_0}$ and fall shorter of $\theta_{n_0}$, and the gap is the sum of two parts that the default condition cannot separate.

\paragraph{P2, norm sensitivity, the primary prediction.} For a second norm $n_1$,
\begin{equation}
\label{eq:p2-group}
G(n_1) - G(n_0) \;=\; \Delta M(n_1) - \Delta M(n_0),
\end{equation}
whatever $\Delta a$ is; and within each person the baseline cancels,
\begin{equation}
\label{eq:p2-within}
c_i(n_1) - c_i(n_0) \;=\; m_i(n_1) - m_i(n_0) \;+\; \varepsilon_i(n_1) - \varepsilon_i(n_0),
\end{equation}
so that a participant's change in effort is predicted, up to noise, by the change in a quantity measured before the interaction. Equation~\eqref{eq:p2-within} is the strongest test the model offers, because it runs on premeasured profiles, person by person, and because the restrictions $M_\gamma$ and $M_\mu$ of \autoref{sec:family} predict a slope of zero for it.

\paragraph{P3, reversal, the strong form.}
\begin{proposition}[The reversal condition]
\label{prop:reversal}
Suppose $G(n_0) > 0$. The gap reverses under $n_1$, $G(n_1) < 0$, exactly when either of the following equivalent conditions holds:
\begin{enumerate}[label=(\alph*), leftmargin=*, nosep]
\item the norm-induced reduction exceeds the initial total gap, $\Delta M(n_0) - \Delta M(n_1) > G(n_0)$;
\item under the new norms the mismatch term favours the minority by more than the baseline gap, $-\Delta M(n_1) > \Delta a$.
\end{enumerate}
\end{proposition}
\begin{proof}
By \eqref{eq:gap}, $G(n_1) = G(n_0) - [\Delta M(n_0) - \Delta M(n_1)]$, which is negative exactly under (a); and $G(n_1) = \Delta a + \Delta M(n_1)$, which is negative exactly under (b).
\end{proof}

\noindent The condition is stricter than the natural reading of ``the norm effect outweighs the baseline''. Take $\Delta a = 3$ and let the mismatch contribution move from $+2$ under $n_0$ to $-2$ under $n_1$. The reduction is $4$, larger than the baseline gap, but smaller than $G(n_0) = 5$: the gap falls from $5$ to $1$ and does not reverse. Reversal needs a reduction greater than $5$, equivalently a new mismatch contribution below $-3$; Figure~\ref{fig:reversal} draws the arithmetic. In the language of aptitude-by-treatment interaction, P2 is an interaction of group with condition, and P3 is the condition for it to be disordinal, the groups' lines crossing between $n_0$ and $n_1$.

\begin{figure}[t]
\centering
\includegraphics[width=0.78\textwidth]{fig_reversal.pdf}
\caption{The reversal condition of \autoref{prop:reversal}. The gap under the new norms falls linearly in the norm-induced reduction; it closes when the reduction equals the whole initial gap $G(n_0)$, not when it equals the baseline gap $\Delta a$. The marked point is the worked case in the text: a mismatch contribution moving from $+2$ to $-2$ reduces the gap by $4$, more than the baseline gap of $3$, and the gap falls from $5$ to $1$ without reversing.}
\label{fig:reversal}
\end{figure}

\section{The model family, and the test that separates its members}
\label{sec:family}

\subsection{One model, four restrictions}

Add to \eqref{eq:cost} the two terms a rival would supply: a condition effect $\gamma_n$ common to all participants, and a mismatch with the partner. For participant $i$ with partner $j$ under norm $n$,
\begin{equation}
\label{eq:family}
c_i(n, j) \;=\; a_i \;+\; \gamma_n \;+\; m_i(n) \;+\; \mu\, d\big(s_i, s_j\big) \;+\; \varepsilon_i(n).
\end{equation}
The models compared are restrictions of \eqref{eq:family}, and each makes a definite prediction for a change of rule set with the partners held fixed (Table~\ref{tab:family}).

\begin{table}[t]
\centering
\small
\caption{The model family: four restrictions of \eqref{eq:family}, and what each predicts when the governing rule set changes and the partners stay the same.}
\label{tab:family}
\begin{tabular}{@{}llp{7.6cm}@{}}
\toprule
Model & Restriction & Prediction under a change of rule set, partners fixed \\
\midrule
$M_\gamma$, general protocol benefit & $\lambda_s = \lambda_k = \mu = 0$ & Effort moves by $\gamma$ for everyone; nothing is redistributed, and nothing is moderated by premeasured mismatch. \\
\addlinespace
$M_\mu$, partner mismatch & $\lambda_s = \lambda_k = 0$ & Beyond $\gamma$, nothing changes, since the partner is unchanged. \\
\addlinespace
$M_\lambda$, norm mismatch & $\mu = 0$ & Effort shifts toward whoever mismatches the governing rules, in proportion to premeasured mismatch (P2), and reverses only under the condition of \autoref{prop:reversal}. \\
\addlinespace
$M_{\mathrm{full}}$ & none & Both the partner and the norm terms. \\
\bottomrule
\end{tabular}
\end{table}

\noindent No tradition is identified with a restriction. An account that emphasizes individual limitations can predict effects of task demands and of accommodations, and so is not $M_\gamma$ or $M_\mu$ by fiat; a double-empathy account can include norm effects, and so is not $M_\mu$ by fiat. The restrictions are compared as models, by what they predict for the data of \autoref{sec:design}.

\subsection{The double-empathy model, read from its text}

The formal model the double-empathy literature has is the distinct communication strategies model \cite{Calderoli2026}, and its characterization here is taken from the paper's own text rather than from its abstract.
\begin{itemize}[leftmargin=*]
\item \emph{The implemented mechanism} is channel preference. Each agent registers the partner's empathy as a weighted average of the partner's verbal and nonverbal output, with the verbal weight the agent's own preference $cp$; the gap between the empathy the agent expects and the empathy it registers raises its defensiveness, which lowers its own output, and the loop can collapse.
\item \emph{What the text allows.} It states that expectations ``may depend on context, relationship history, communication norms'', and holds expected empathy at one value for both agents in its simulations. It reports that ``the trajectory to collapse shows greater experiential asymmetry'', and predicts, as a falsification criterion, that ``the partner with higher $cp$ should report lower perceived empathy''. Its limitations section says that ``the costs are therefore bilateral but need not be equal or arise from the same source''.
\item \emph{What it has no parameter for.} An imposed or governing norm. Its intervention prediction concerns training that increases verbal empathy expression under stress.
\end{itemize}
It would therefore be wrong to say that this model predicts equal costs, or invariance under a change of norms; it is silent on norms as implemented, and open to them in its text. The distinction between it and the norm-setting model has to be shown in data, and the data can show it only under assumptions, which are stated next.

\paragraph{The mediation analysis, and what it can establish.} Suppose the double-empathy mechanism is extended so that expected empathy depends on the governing norm. A change of rule set then acts on effort \emph{through} the gap between the empathy a participant expects and the empathy the participant perceives: the gap moves, and effort moves with it. Under $M_\lambda$ effort moves with premeasured profile-to-norm mismatch whether or not the gap moves. Three things follow for the design, and together they are less than a test.
\begin{enumerate}[label=(\roman*), leftmargin=*]
\item \emph{The mediator is the gap, not perceived empathy.} If under rule set A a participant expects $8$ and perceives $6$, and under rule set B expects $6$ and perceives $6$, the gap falls from $2$ to $0$ and the whole change in effort may run through it while perceived empathy has not moved; an analysis that controlled perceived empathy alone would count that change as a norm effect beside the mechanism. The design therefore measures both, expected partner empathy before each condition's task and perceived partner empathy after it, and the mediator is their difference.
\item \emph{Controlling the mediator does not identify the direct effect by design.} The rule set is randomized; the gap is not. The coefficient of premeasured mismatch in the within-person regression, with the change in the gap controlled, is the direct effect only if the gap's change is unconfounded with the effort change given the rule set and the premeasured profiles, the sequential ignorability of causal mediation analysis \cite{ImaiKeeleTingley2010, ImaiKeeleYamamoto2010}, and measurement error in the empathy reports, common causes of the reports and of effort, and feedback from effort to the judgement of the partner each violate it. The analysis is therefore run as an assumption-dependent investigation: direct and indirect effects estimated under sequential ignorability, the sensitivity analysis of Imai, Keele and Yamamoto \cite{ImaiKeeleYamamoto2010} reporting the mediator--outcome error correlation at which each would vanish, and the reliabilities of the empathy reports, estimated in the pilot, used to correct the attenuation.
\item \emph{The two accounts are not exclusive.} Equation~\eqref{eq:cost} does not say by what route mismatch produces effort, and a positive $\lambda$ may describe an effect that runs in part through the gap. A direct effect that shrinks when the gap is controlled is consistent with the cost model; no categorical prediction, zero, the full slope or something between, is made for the residual coefficient; and a direct effect that vanishes at a stated precision would show that the first study cannot find a norm effect beside the gap, not that none exists.
\end{enumerate}
What the first study can establish is whether the mismatch effect is present with the gap controlled, and how much of it the gap carries under the stated assumptions and across the stated sensitivity range. What would separate the mechanisms is a study that manipulates expected empathy directly, by disclosure or by instruction, which is deferred (\autoref{sec:design}).

\subsection{The established setting}

Context-dependent advantages that reverse across conditions are aptitude-by-treatment interactions \cite{Cronbach1957, CronbachSnow1977}: the person-by-condition interaction that Cronbach urged the two disciplines of psychology to study together, and whose disordinal form is the case in which no condition is best for everyone. The norm-setting model is an instance, with two specifications that give it content. The aptitude is a premeasured profile, not a trait inferred from the outcome, and the treatment is a communication norm imposed on both partners of a dyad at once, so that the interaction of interest is between the norm and the \emph{distribution} of cost within the pair. Within that setting the paper's claim is the one of \autoref{sec:intro}: that the distribution is estimable, and separately from the partner effect under an assumption stated in \autoref{sec:simulation}, by the design that follows.

\section{Study 1: a discriminating design}
\label{sec:design}

The design manipulates the governing rules, holds fixed what can be held fixed of everything else, and measures what cannot. A comparison of ``default conversation'' with ``an explicit protocol'' would change norms, structure, predictability, memory demands and perhaps modality at once, and an improvement under the protocol could be any of them. The study therefore compares two explicit rule sets.

\paragraph{Participants.} Mixed dyads, one autistic and one non-autistic adult, with diagnosis by report and the autistic participants' self-identification recorded. Same-neurotype dyads are a later study (see ``Deferred'' below).

\paragraph{Before any interaction.}
\begin{enumerate}[label=(\roman*), leftmargin=*]
\item Each participant's style profile $s_i$ and competence profile $k_i$ are measured with instruments fixed in the preregistration \cite{Nosek2018}: for style, self-report and a brief performance measure of directness, literalness and turn-yielding preference; for competence, a task of implicature comprehension and a task of following explicit structure, each scored on the scale the requirement is written in.
\item Independent raters, blind to the outcomes, code each rule set's presupposed style $s^{*}_n$ and competence requirement $\theta_n$ from its text, on the same scales.
\item The distance $d$, the weights $w_n$, the outcome measures and the analysis are fixed in the preregistration. Profiles, metric and norm codings chosen after the outcomes are seen would make the model too flexible to fail.
\item The common metric is established, not assumed. Independent coding protects the norm locations from being chosen to fit the outcomes; it does not make a participant's preference score and a coded norm location commensurable. The pilot therefore codes the rule sets on the anchored scales the profile instruments report on, infers the competence thresholds from the rule text on the scale the competence tasks are scored in, and checks the calibration against behaviour: participants' adherence to each rule set and their reported difficulty with it should order with the coded distance, and where they do not the scales are revised before the main study.
\end{enumerate}

\paragraph{The manipulation.} Two explicit rule sets differ in the communicator they presuppose and are matched on specified incidental features: length, number of rules, channel, time, and the structure of the task they govern. They are not matched in everything, and the model does not ask them to be. The competence term says that the rule sets draw on different competences; inferring a hint and yielding a turn on an implicit cue make demands that stating a request and yielding in words do not, and those are the differences the model is about. What remains, the learning burden of each rule set and its automaticity after practice, is measured rather than assumed equal: training time to criterion, practice errors and reported effort during practice are recorded for each rule set, enter the analysis as covariates, and serve as the check on the matching, which is treated as a claim the pilot supports or revises.
\begin{itemize}[leftmargin=*]
\item \emph{Rule set A} encodes indirect, implicature-rich conventions: requests made by hinting, an opening exchange of small talk, turns yielded by implicit cues.
\item \emph{Rule set B} encodes direct, literal conventions: requests stated outright, no obligatory small talk, turns yielded in words.
\end{itemize}
Both are practised to a criterion before the task; adherence to the rules, coded from recordings, is the manipulation check. Which conventions rule set B should encode, and from what sources, is a decision for the preregistration, informed by autistic participants' own accounts of their communication and by the published descriptions of autistic communication preferences; the design does not presume that rule set B is the autistic participants' preferred style, it measures where each participant stands relative to each rule set.

\paragraph{Within dyads.} Each pair works under both rule sets, in counterbalanced order, on two parallel versions of a collaborative task, also counterbalanced. The order is assigned to dyads at random and balanced; a period effect is estimated as the order-by-condition term; differential carryover, which a two-period crossover cannot separate from the order effect \cite{Senn2002}, is guarded against by the parallel task versions and a washout between conditions and is assessed by the adherence coding of the second condition, in which the first rule set's conventions would show if they persisted. The room, lighting and noise are identical across conditions; fatigue is measured after each condition. The pair is its own control: the baseline $a_i$, the partner term $\mu\, d(s_i, s_j)$ and the dyad's shared conditions cancel in the within-person difference \eqref{eq:p2-within}.

\paragraph{Outcomes.} Effort per participant and condition, primary, by two measures fixed in advance: the mental-demand and effort items of the NASA Task Load Index \cite{HartStaveland1988}, reworded to the conversation just held (how mentally demanding the conversation was; how hard the participant had to work to follow the rules and accomplish the task) on the index's scale, and one behavioural index, the rate of repair initiations and restarts per minute of the participant's speech, coded from the recording by raters blind to condition. Low reported effort can mean easier coordination or lower engagement, so effort is interpreted with task performance and adherence beside it: a condition in which effort, performance and adherence all fall is read as disengagement, not as ease. Task performance and rapport, secondary; expected and perceived partner empathy per condition, whose difference is the mediator; rule adherence, the manipulation check; training time, practice errors and practice effort per rule set, the matching check.

\paragraph{Predictions.}
\begin{itemize}[leftmargin=*]
\item \emph{Primary (P2).} On effort, the rule set interacts with premeasured mismatch within persons: the change in effort from A to B is predicted by the change in $m_i$, with slope one in the model's units and positive in any units. The group-level redistribution of effort is this prediction's aggregate.
\item \emph{$M_\gamma$.} A main effect of rule set only, with no moderation by premeasured mismatch.
\item \emph{$M_\mu$.} No change beyond the main effect, the partner being unchanged.
\item \emph{A norm-dependent double-empathy mechanism.} A change carried by the expected--perceived empathy gap. What the first study can say about it is stated in \autoref{sec:family}; it is not a prediction the primary analysis tests.
\end{itemize}

\paragraph{Analysis.} The primary analysis is the within-person difference regression of \eqref{eq:p2-within}, the change in effort on the change in premeasured style mismatch and competence shortfall, with standard errors clustered by dyad. Its primary test is the joint Wald test of the two mismatch coefficients, $\lambda_s$ and $\lambda_k$, at the $5\%$ level, so that primary success is evidence that premeasured mismatch moderates the change in effort, through either term or both; the power calculations of \autoref{sec:simulation} are for this test. Evidence for each coefficient separately is secondary, tested one-sided in the predicted direction with the familywise rate held at $5\%$. Levels are analysed in the actor--partner interdependence model \cite{Kenny2006}, participants nested in dyads and conditions within participants, which identifies $\mu$ from between-dyad variation under the assumption stated in \autoref{sec:simulation}. The four models are compared by their Bayesian information criterion on the exact likelihood of Appendix~\ref{app:estimators}, and the mediation analysis is as stated in \autoref{sec:family}. Appendix~\ref{app:prereg} lists what the preregistration fixes.

\paragraph{What would count against the model.} Stated in advance, with the preregistration, and each with the precision it needs, since a nonsignificant estimate establishes no absence:
\begin{enumerate}[label=(\roman*), leftmargin=*, nosep]
\item a manipulation that changes adherence but moves no one's effort, the condition effect and both mismatch coefficients bounded within their smallest effects of interest;
\item a change of effort that premeasured mismatch does not moderate: the joint test not rejecting, and an equivalence test \cite{Lakens2017, Lakens2018} rejecting, for each coefficient, every value beyond its smallest effect of interest, which the preregistration sets at a quarter of the simulated values, $0.15$ for $\lambda_s$ and $0.20$ for $\lambda_k$ in the model's units, the sizes at which the simulation's power is reported, or at the pilot's estimate where there is one; a nonsignificant moderation whose interval still covers the smallest effect of interest is inconclusive, not a disconfirmation;
\item a direct effect of premeasured mismatch that is bounded within the smallest effect of interest once the change in the empathy gap is controlled, with the sensitivity analysis showing the bound robust across the range of mediator--outcome confounding the pilot's reliabilities make plausible.
\end{enumerate}
The first two falsify the norm-setting model as stated. The third would show that the first study cannot find a norm effect beside the double-empathy mechanism, and would send the question to the study that manipulates expectations.

\paragraph{Deferred.} Composition, all-autistic and all-non-autistic dyads beside mixed ones, and disclosure of diagnosis to the partner are each a manipulation of their own, and each goes to its own study. Composition tests whether the partner term $\mu$ and the norm terms interact, and whether the same-neurotype advantage the first diffusion-chain study found \cite{Crompton2020}, and its replication did not \cite{Crompton2025}, appears under a change of governing norm; disclosure manipulates expectations directly and is the study the mediation analysis of \autoref{sec:family} defers to. A third design, each participant interacting with several partners in a round robin \cite{KennyLaVoie1984}, identifies the partner term within persons rather than under the between-dyad assumption of \autoref{sec:simulation}, and is the stronger design for $\mu$. A first study that crossed norms with composition, disclosure or repeated partners would need several times the dyads and would confound its own manipulation checks.

\section{The simulation}
\label{sec:simulation}

Before any data, the design was run on simulated participants to settle five questions: whether the analysis recovers the parameters of \eqref{eq:family} at the planned size; whether it selects the generating model among the four; whether a model fitted to two rule sets predicts a third, for participants it has not seen; what the reversal arithmetic looks like at realistic profiles; and, since each of these answers is conditional on the generating model, which of them survive departures from it. Showing that the equations generate the intended pattern would establish much less than this, since any model generates its own pattern. The simulation is the supplementary script \texttt{norm\_cost\_sim.py}, whose complete output accompanies the paper; its settings are in Table~\ref{tab:sim-settings}.

\begin{table}[t]
\centering
\small
\caption{The simulation's settings. Style and competence profiles are two-dimensional, drawn around group means that overlap; rule sets A and B presuppose opposite styles and draw on different competences; rule set C is withheld from fitting. Effort is in standard-deviation units of the baseline.}
\label{tab:sim-settings}
\begin{tabular}{@{}lp{6.1cm}p{4.9cm}@{}}
\toprule
Quantity & Setting & Role \\
\midrule
Dyads & $60$ mixed, each under A and B & the planned size; $20$ to $120$ in the power study \\
Style profiles $s_i$ & $\mathcal{N}(\pm(0.8, 0.8),\, 0.7^2 I)$ & autistic $+$, non-autistic $-$, on (directness, literalness) \\
Competence profiles $k_i$ & $\mathcal{N}(\pm(-0.3, 0.3),\, I)$ & autistic $+$, non-autistic $-$, on (implicature reading, explicit structure) \\
Rule set A & $s^{*} = (-1, -1)$, $\theta = (1, 0)$, $w = (1, 0)$ & indirect; requires implicature reading \\
Rule set B & $s^{*} = (+1, +1)$, $\theta = (0, 1)$, $w = (0, 1)$ & direct; requires following explicit structure \\
Rule set C & $s^{*} = (+1, -1)$, $\theta = (0.5, 0.5)$, $w = (0.5, 0.5)$ & withheld; direct but non-literal \\
$\lambda_s, \lambda_k, \mu, \gamma_B$ & $0.6$, $0.8$, $0.4$, $-0.3$ & the generating values; $\gamma_A = 0$ \\
$\Delta a$ & $0.5$ & the autistic group's baseline shift \\
$\sigma_a, \sigma_u, \sigma_\varepsilon$ & $1.0$, $0.4$, $0.8$ & baseline, dyad and residual standard deviations \\
\bottomrule
\end{tabular}
\end{table}

\paragraph{Estimation.} The simulation generates effort from \eqref{eq:family} with one addition, a dyad effect shared by the two partners, which the within-person difference cancels. The balanced crossover decomposes exactly into within-person differences, which identify $\gamma$ and the two $\lambda$'s with the baseline, the dyad effect and the partner term cancelled, and person means, which identify $\mu$ beside a group indicator carrying the baseline gap. The second identification is conditional in a way the first is not. With one partner per participant, $a_i' = a_i + \mu\, d(s_i, s_j)$ with $\mu' = 0$ gives the same observations as $a_i$ with $\mu$, so the person means identify $\mu$ only under a restriction on how baselines relate to partner distance: the simulation generates baselines mean-independent of partner distance given the group and the measured profiles, and that is the assumption the real study's between-dyad estimate needs, one the actor--partner model states as a random-effects assumption and the design does not supply. People with unusual profiles may have both more distant partners and higher baselines for reasons outside the measured mismatch, and the robustness study below shows what that does. Each part is fitted by least squares with dyad-clustered standard errors, the primary test is the joint Wald test of the two $\lambda$'s on the clustered covariance, and the model comparison uses the exact likelihood of Appendix~\ref{app:estimators}, in which the dyad effect's covariance between partners is carried by a rotation of the two partners' means.

\paragraph{Parameter recovery.} Over $1{,}000$ simulated studies of $60$ dyads, every parameter is recovered without bias and with coverage close to nominal (Table~\ref{tab:recovery}, \autoref{fig:recovery}). The difference equation identifies three parameters within persons, and not equally well: the standard deviations of $\hat\lambda_s$ and $\hat\lambda_k$ are a tenth of their values and the joint test rejects in every replication, while $\gamma$, a difference of condition means against the residual noise, has power $81\%$. The partner term $\mu$, identified between dyads from the variation of partner distance against the baseline variance, has power $77\%$. The power study (\autoref{fig:power}) puts the joint test at ceiling from $20$ dyads at the generating effect sizes, which says only that those sizes are large. At a quarter of them, $\lambda_s = 0.15$ and $\lambda_k = 0.20$, the joint test has power $66\%$ at $20$ dyads, $93\%$ at $40$ and $99\%$ at $60$, and with the profiles and the codings measured with error $60\%$, $87\%$ and $95\%$: the planned $60$ dyads hold $90\%$ power down to a quarter of the simulated mismatch effects under measurement error, which is the figure the preregistration should quote. $\gamma$ passes $80\%$ at $60$ dyads, $\mu$ $90\%$ at $90$ and $96\%$ at $120$. When the truth is $M_\gamma$ the false-positive rate is $4.9\%$ for the joint test and $5.1\%$ and $5.3\%$ for the separate tests, but $9.8\%$ for ``either coefficient'', which is why the joint test is primary and the separate tests are secondary.

\begin{table}[t]
\centering
\small
\caption{Parameter recovery over $1{,}000$ simulated studies of $60$ mixed dyads, the full model fitted: the generating value, the mean estimate, its standard deviation, the coverage of the dyad-clustered $95\%$ interval, and the power of the two-sided $5\%$ test against zero. The primary joint test of $(\lambda_s, \lambda_k)$ has power $100\%$ at these effects and a false-positive rate of $4.9\%$ when the truth is $M_\gamma$.}
\label{tab:recovery}
\begin{tabular}{@{}lrrrrr@{}}
\toprule
Parameter & Truth & Mean estimate & SD & $95\%$ coverage & Power \\
\midrule
$\gamma$ (condition effect) & $-0.300$ & $-0.296$ & $0.102$ & $95.1\%$ & $81.1\%$ \\
$\lambda_s$ (style mismatch) & $0.600$ & $0.601$ & $0.057$ & $94.4\%$ & $100\%$ \\
$\lambda_k$ (competence shortfall) & $0.800$ & $0.802$ & $0.085$ & $93.5\%$ & $100\%$ \\
$\mu$ (partner mismatch) & $0.400$ & $0.402$ & $0.148$ & $94.8\%$ & $77.4\%$ \\
\bottomrule
\end{tabular}
\end{table}

\begin{figure}[t]
\centering
\includegraphics[width=\textwidth]{fig_recovery.pdf}
\caption{Sampling distributions of the four estimates over $1{,}000$ simulated studies of $60$ dyads, the full model fitted. The vertical line marks the generating value. The mismatch parameters $\lambda_s$ and $\lambda_k$ are estimated with a precision that $\gamma$, identified within persons as a difference of condition means, and $\mu$, identified between dyads, do not reach at this size.}
\label{fig:recovery}
\end{figure}

\begin{figure}[t]
\centering
\includegraphics[width=0.8\textwidth]{fig_power.pdf}
\caption{Power against the number of mixed dyads, $400$ simulated studies per point. Solid lines: the primary joint test of $(\lambda_s, \lambda_k)$ at the generating effects, and the two-sided tests of the condition effect $\gamma$ and the partner term $\mu$. Dashed: the joint test at a quarter of the mismatch effects ($\lambda_s = 0.15$, $\lambda_k = 0.20$); dotted: the same with the profiles and the rule-set codings measured with error. The joint test is at ceiling at the generating effects; at a quarter of them it passes $90\%$ at $40$ dyads, and at $60$ under measurement error; $\gamma$ passes $80\%$ at $60$ dyads and $\mu$ $90\%$ at $90$.}
\label{fig:power}
\end{figure}

\paragraph{Model recovery.} With $500$ studies generated from each model and the four fitted to each, the Bayesian information criterion on the exact likelihood selects the generating model in $98\%$ of the $M_\gamma$ studies, $76\%$ of $M_\mu$, $96\%$ of $M_\lambda$ and $80\%$ of $M_{\mathrm{full}}$ (\autoref{fig:confusion}, left). The two large errors are of one kind: the partner term is dropped, $M_\mu$ mistaken for $M_\gamma$ in $24\%$ of its studies and $M_{\mathrm{full}}$ for $M_\lambda$ in $20\%$, because $\mu$ is the parameter the design identifies least well. The norm terms are almost never mistaken: no study generated with norm mismatch selected a model without it, and $5$ of the $1{,}000$ generated without it selected a model with it. A predictive log score with whole dyads held out, six folds, is reported beside the criterion (right) as the comparison that needs no assumption about the effective sample size. It recovers $M_{\mathrm{full}}$ as well as the criterion does and never drops the norm terms when they are present, but it over-selects when the truth is small, choosing a model with the norm terms in $32\%$ of the $M_\gamma$ studies and $32\%$ of the $M_\mu$ studies: the known behaviour of such scores, since leave-one-out is asymptotically equivalent to Akaike's criterion \cite{Stone1977} and cross-validation that holds out a fixed fraction of the data, as six folds do, does not select consistently among nested models \cite{Shao1993}. The preregistered criterion is therefore the Bayesian one, and whether the norm terms are present at all is settled by the primary test and not by the model comparison.

\begin{figure}[t]
\centering
\includegraphics[width=0.92\textwidth]{fig_confusion.pdf}
\caption{Model recovery: the share of simulated studies ($500$ per generating model, $60$ dyads) in which each of the four models was selected, by the Bayesian information criterion on the exact likelihood (left) and by the dyad-held-out predictive log score (right). Under the criterion the off-diagonal mass lies almost entirely in dropping the partner term; the predictive score over-selects the larger models when the truth is small.}
\label{fig:confusion}
\end{figure}

\paragraph{The P2 signature.} \autoref{fig:signature} shows one simulated study under each of two truths. Under norm mismatch the within-person change in effort from A to B falls on a line through the change in premeasured mismatch, with the autistic participants, whose profiles lie nearer rule set B, moving down and the non-autistic participants moving up; under a general protocol benefit the same plot is flat, both groups moving together by $\gamma$. This is the figure the real study will produce, and the two panels are what its two main outcomes look like. In the figure the two premeasured mismatches are combined into one score with the generating weights, fixed before the data. In the real study the weights that combine them are either fixed in advance from the pilot or estimated on dyads other than those plotted: a plot whose weights are fitted to the effort data it displays shows the fit and not a prediction, whatever its slope.

\begin{figure}[t]
\centering
\includegraphics[width=\textwidth]{fig_signature.pdf}
\caption{The P2 signature in one simulated study of $60$ dyads under each of two truths. Each point is a participant: the change in effort from rule set A to rule set B against the change in premeasured mismatch, combined with weights fixed before the data, in the model's cost units. Under norm mismatch (left) the change lies on a line of slope near one and the groups move in opposite directions; under a general protocol benefit (right) the slope is near zero and the groups move together.}
\label{fig:signature}
\end{figure}

\paragraph{The withheld rule set.} The mismatch coefficients were fitted to rule sets A and B on five sixths of the dyads and used to predict, for the participants of the held-out sixth, the change in effort from A to the withheld rule set C from their premeasured profiles alone, over six folds. What is predicted is the relative individual change: A and B cannot identify the condition effect $\gamma_C$, so the predicted and the observed changes are each centred within the held-out set, and the exercise tests whether the model says who will find C harder than A, and by how much relative to the others, not the absolute effect of C. With the truth $M_\lambda$ the root-mean-square error is $1.10$ (standard deviation $0.07$ over $500$ studies), against $1.47$ for the rivals' prediction of no individual differences; the irreducible error of a centred within-person change is $\sqrt{2}\,\sigma_\varepsilon \sqrt{1 - 1/20} = 1.10$ at this fold size, so the norm models predict the relative effects of a rule set they have not seen, for participants they have not seen, to the noise floor. The same holds with the truth $M_{\mathrm{full}}$ ($1.11$ against $1.46$). Holding out dyads removes the dependence that reusing the same participants' A observations in training and in evaluation would create; what remains is that the observed change $y_{i,\mathrm{C}} - y_{i,\mathrm{A}}$ carries the noise of two observations, which is what sets the floor at $\sqrt{2}\,\sigma_\varepsilon$, and a replicated baseline measurement would lower it; the real study's third rule set, if one is run, is where that replication belongs.

\paragraph{The reversal arithmetic.} At the generating values, on a population of $20{,}000$ dyads, the autistic group's mean style mismatch is $2.64$ under rule set A against the non-autistic group's $0.91$, and $0.92$ against $2.65$ under rule set B; the mean competence shortfalls are $1.35$ against $0.83$ and $0.85$ against $1.34$. The mismatch gap $\Delta M$ is $+1.45$ under A and $-1.44$ under B, so that with the baseline gap of $0.50$ the total gap $G$ is $+1.95$ under A and $-0.94$ under B: the norm-induced reduction, $2.89$, exceeds $G(\mathrm{A}) = 1.95$, and the gap reverses, as \autoref{prop:reversal} requires. Under the withheld rule set C, which presupposes a style equidistant from both groups' means, $\Delta M$ is $-0.003$ and the gap is the baseline alone, $0.50$. The example shows the three regimes the proposition distinguishes in one population: a gap set by the norms, a gap reversed by them, and a gap the norms leave to the baseline.

\paragraph{Robustness.} Every result above is conditional on the generating model, and the real study will not satisfy it. Table~\ref{tab:robustness} repeats the recovery under five departures from it, singly and together, at $60$ dyads. Three conclusions survive all of them: the primary joint test keeps its power, the Bayesian information criterion keeps selecting a model with the norm terms, and the within-person estimates $\hat\lambda_s$ and $\hat\lambda_k$ stay near their values. Two do not. Measurement error in the profiles and the codings attenuates $\hat\lambda_k$ by an eighth and, more consequentially, spoils the intervals: a rule set's coding error is one draw shared by every participant in the study, so it does not average out within a study, and the coverage of $\hat\gamma$'s interval, which absorbs it, falls to $55\%$, that of $\hat\lambda_k$'s to $74\%$. The real study's intervals are therefore honest only if the coding error is small or is propagated, by coding each rule set with several raters and carrying their disagreement into the standard errors, and the pilot's inter-rater reliability decides which. And baselines correlated with partner distance bias $\hat\mu$ by more than its value, with coverage $7\%$, as the identification argument of the estimation paragraph says they must, while leaving the within-person estimates untouched: the partner parameter is the one the design does not protect, and its estimate in the real study is to be read under the stated assumption or not at all.

\begin{table}[t]
\centering
\small
\caption{Robustness to departures from the generating model: $60$ dyads, $300$ simulated studies per row, the full model fitted. Each cell gives the bias of the estimate and, in parentheses, the coverage of its nominal $95\%$ interval in percent; the last two columns give the power of the primary joint test and the share of studies in which the Bayesian information criterion selected a model with the norm terms, in percent. The third and the last rows generate the mismatch effects at half their values ($\lambda_s = 0.30$, $\lambda_k = 0.40$); measurement error adds noise of standard deviation $0.5$ to each profile coordinate and $0.3$ to each coded style and threshold coordinate, one draw per study; the correlated baselines add $0.5$ per unit of partner distance and $0.4$ per unit of profile extremity; the unequal variances are $0.6$ under A and $1.0$ under B; the correlated predictors carry $0.6$ of the centred style into competence.}
\label{tab:robustness}
\begin{tabular}{@{}>{\raggedright\arraybackslash}p{4.7cm}cccccc@{}}
\toprule
Departure & $\gamma$ & $\lambda_s$ & $\lambda_k$ & $\mu$ & Joint & BIC \\
\midrule
As generated & $-0.010$ (95) & $-0.002$ (97) & $+0.002$ (96) & $+0.001$ (92) & $100$ & $100$ \\
\addlinespace
Profiles and codings measured with error & $-0.006$ (55) & $+0.017$ (84) & $-0.095$ (74) & $+0.070$ (91) & $100$ & $100$ \\
\addlinespace
Mismatch effects halved & $+0.001$ (94) & $-0.005$ (94) & $+0.009$ (94) & $+0.020$ (93) & $100$ & $100$ \\
\addlinespace
Baselines correlated with profiles and partner distance & $-0.003$ (95) & $-0.001$ (94) & $+0.003$ (95) & $+0.556$ (7) & $100$ & $100$ \\
\addlinespace
Unequal residual variances across conditions & $+0.007$ (97) & $+0.004$ (92) & $+0.003$ (95) & $+0.001$ (93) & $100$ & $100$ \\
\addlinespace
Style and competence correlated within person & $-0.003$ (95) & $+0.000$ (91) & $-0.009$ (92) & $+0.005$ (94) & $100$ & $100$ \\
\addlinespace
All five together, effects halved & $+0.012$ (80) & $+0.004$ (91) & $-0.055$ (89) & $+0.665$ (2) & $100$ & $100$ \\
\bottomrule
\end{tabular}
\end{table}

\paragraph{What the simulation does not show.} It does not show that the effect sizes are realistic; they are assumptions, chosen so that the groups' style distributions overlap and the baseline gap is a third of the norm-induced one, and the quarter-effect power curves are the nearest the simulation comes to a guard against them. It does not show that style and competence can be measured on scales commensurate with the codings of the rule sets, which is the measurement problem the pilot has to solve. And it does not model the empathy gap, whose analysis depends on assumptions that a simulation building them in could not test.

\section{What existing results constrain}
\label{sec:constraints}

Three bodies of results bear on the design, and each is read from its primary source.

\paragraph{Information transfer in chains.} The first diffusion-chain study found that information passed along chains of autistic adults, and along chains of non-autistic adults, more effectively than along mixed chains \cite{Crompton2020}. The registered replication with $311$ participants did not find that difference: there was no main effect of chain type on information transfer, for fictional or for factual stories, and an unpredicted interaction between autistic chain type and position for fictional stories, performance falling more steeply along autistic chains than non-autistic ones \cite{Crompton2025}. Rapport, self-rated and analysed by role, was higher for teachers in non-autistic chains than in mixed and in autistic chains, and for learners in non-autistic chains than in mixed chains only; telling participants their chain's neurotype raised rapport significantly for learners and not for teachers. For this paper the constraint is twofold: information transfer is not where the cost appeared in that task, and the rapport effects depend on role, which is why effort rather than rapport is the primary outcome here and why rapport gets a model of its own.

\paragraph{Rapport in groups.} In $36$ groups of four, all-autistic, all-non-autistic and three-to-one majorities, there was no overall group-type difference in composite rapport \cite{Foster2025}. All-autistic groups rated their interaction more enjoyable and friendlier than both kinds of mixed group; autistic participants' ease and enjoyment fell as non-autistic members were added, the composite interaction not reaching significance, while non-autistic participants' composite rapport stayed stable; the item-level results are uncorrected for multiple comparisons and are called preliminary by their authors. Neither this study nor the chain study shows that majority-composition settings favour non-autistic rapport in general. The group study did not manipulate norms, so any reading of its item-level pattern through the cost model, the autistic participants' falling ease as the session's conventions shift toward the majority's, is a hypothesis for Study 1 and not a result.

\paragraph{Perception and the double-empathy model.} Non-autistic observers form less favourable first impressions of autistic adults from brief samples, on the basis of presentation rather than content \cite{Sasson2017}, which is one reason the cost of a mixed interaction cannot be read off its outcome alone. The distinct communication strategies model \cite{Calderoli2026} is characterized in \autoref{sec:family}; it supplies the paper's strongest formal rival and the reason the empathy gap is measured.

\section{Outside this paper}
\label{sec:extensions}

Three things the source raises are extensions, and the paper says of each what it would take.

\paragraph{Branch geometry.} The Branch Principle is a hypothesis about the structure of cognitive profiles: at matched resolution, non-autistic profiles form one type, or few, and autistic profiles form several, in the strong form about as distant from each other as from the non-autistic type. Its test fits the same latent-structure models to both groups on the same cognitive-profile measures, a mixture model, an additive tree \cite{SattathTversky1977} and a dimensional model, and compares them by likelihood; the recent decomposition of autism's phenotypic heterogeneity into classes with distinct genetic programs \cite{Litman2025} is evidence about clinical phenotype at a different grain, not about the cognitive profiles the principle concerns. The geometry is independent of the cost model and is not support for mutual advantage: equal distances do not imply crossed advantages. With $r = (0, 0)$, $b = (1, 1)$ and $c = (1, -1)$, both $b$ and $c$ lie at distance $\sqrt{2}$ from $r$, yet $b$ dominates $r$ while $c$ and $r$ cross. Linking the geometry to \autoref{prop:reversibility}'s premise needs an argument of its own, and an asymmetric notion of similarity \cite{Tversky1977} may be the right instrument for it.

\paragraph{Linguistic masking.} The source names a phenomenon a layer below behavioural camouflaging \cite{Hull2017, Hull2019, Pearson2021}: a speaker who learns a cognitive term's public extension case by case, over a private intension that differs from the dominant group's, so that the difference stays hidden on familiar cases and shows only when the term is extended to new ones. Its signature is a generalization gradient, agreement with the dominant group on familiar instances and divergence on novel ones. The gradient is necessary for the phenomenon and not unique to it, since differences in concept knowledge, familiarity, confidence or learning history produce the same pattern, and response strategies shift with uncertainty. What would tie the gradient to the mechanism is a private organization of the instances measured without the term, by similarity sorting, predicting the divergence on novel instances over and above familiarity and confidence; a larger gradient for terms learned from usage than for terms learned from definitions; and an experimental model in which a private organization of a novel domain is induced before a public term is taught over a different one. The question of which instruments the phenomenon would and would not reach is open: an eight-item alexithymia score shows practically ignorable differential item functioning between autistic and non-autistic adults \cite{WilliamsGotham2021}, so the claim cannot be that every self-report instrument distorts, only that particular wording features might. This is a second paper.

\paragraph{Shared absence.} A further hypothesis, that non-reference branches share the absence of the reference's conventions and so find each other's interactions cheaper even where their own preferences differ, would separate rapport from the prediction of preferences. Its added premises are stated when it is tested, against familiarity, stigma, preference variability and disclosure, and it is not an advance prediction of anything reported here.

\section{Discussion}
\label{sec:discussion}

\paragraph{What the model adds.} The received descriptions of a mixed interaction's difficulty place it in a person or in a pair. The cost model places part of it in the rules, and makes the placement measurable: the share of the group gap that moves when the rules move, person by person, from profiles measured in advance. That share is what Study 1 estimates, and the model family is what keeps the estimate honest, since a general benefit of explicit rules and a mismatch with the partner are both allowed to carry the data if they can.

\paragraph{The conditional status of the conclusions.} Nothing above establishes that effort tracks mismatch; the paper is built so that the question can be settled. The exact results, reference-relativity and the reversal condition, hold of the definitions and the model; the cost model is a hypothesis; the Branch Principle is an extension with its own test; balance is a conjecture. The reversal condition in particular should be read as a caution to the design: a study that finds the gap reduced by more than the baseline gap has not found a reversal, and a study powered to find the reduction is not thereby powered to find the crossing.

\paragraph{Limitations the design carries.}
\begin{itemize}[leftmargin=*]
\item \emph{Measurement.} The model's content is only as good as the commensurability of the profiles and the norm codings, and the choice of the distance $d$ and of the weights $w_n$. All are fixed in advance, which protects the test from being fitted to its outcome, but not from a poor choice: mismeasured profiles attenuate the mismatch slopes, a coding error shared by every participant biases the condition effect and no interval widens to show it, and a metric that places profiles and codings on scales that are not in fact common can manufacture an association as easily as hide one. The pilot's calibration is the remedy, and the robustness study says what is at stake.
\item \emph{Two rule sets.} Two norms identify the norm terms but not their functional form; a third rule set, as in the simulation's withheld condition, is the natural extension and the strongest test.
\item \emph{Dyads.} The within-person test is clean, but the partner term is identified between dyads, needs $90$ or more pairs at the simulated effect size, and rests on an assumption about the baselines that the design cannot enforce (\autoref{sec:simulation}); a study designed for the primary prediction alone is underpowered for the full family, and a round-robin design is the one that identifies the partner term within persons.
\item \emph{Scope.} Both conditions impose explicit instructions, so what the first study measures is adaptation to experimental protocols. Extending its result to the norms of ordinary interaction needs evidence that the protocols represent them, which a comparison of the pilot's adherence codings with recordings of unconstrained conversation can begin to supply.
\item \emph{Generality.} The groups studied are autistic and non-autistic adults, since that is where the question arose and where the literature is; the model is stated for any two groups whose communicative profiles differ, and its application elsewhere is a hypothesis in each case.
\item \emph{Participation.} The rule sets, the profile instruments and the interpretation of effort should be developed with autistic researchers and participants, not only about them; the design's choices that are left open are left open for that reason.
\end{itemize}

\paragraph{What would advance the account most.} Not a further qualification, but the study itself, preregistered, with its third rule set; and, for the extensions, a cognitive-profile data set at matched resolution on which the branch geometry can be fitted rather than argued.

\paragraph{A note on provenance.} The Branch Principle and the Principle of Symmetry were stated by the author in a short video recorded on June 8, 2024, at 2:39 PM ET, and written into the larger work's section on cognitive style and convention \cite{Aethus}, where the observation that ``a limitation does not toggle'' appears in a separate section on the modal structure of language. The reading of a diagnostic label as a partial information state is the larger work's throughout. The restatement of the Principle of Symmetry as reversibility at a measurement margin, the cost model, the model family, the mediation analysis, the design and the simulation were developed in October 2026; the source passages are quoted as written, and the formal apparatus is the development's.

\printbibliography

\appendix
\section{The simulation's estimators}
\label{app:estimators}

For participant $i$ in dyad $d(i)$ and rule sets $\mathrm{A}$ and $\mathrm{B}$, write $y_{i,n}$ for the effort under \eqref{eq:family}. The within-person difference and the person mean are
\begin{align}
y_{i,\mathrm{B}} - y_{i,\mathrm{A}} &= \gamma + \lambda_s\, \Delta d_i + \lambda_k\, \Delta \mathrm{sh}_i + \big(\varepsilon_{i,\mathrm{B}} - \varepsilon_{i,\mathrm{A}}\big), \label{eq:within}\\
\tfrac12\big(y_{i,\mathrm{B}} + y_{i,\mathrm{A}}\big) &= \alpha + \Delta a\, g_i + \tilde a_i + u_{d(i)} + \lambda_s\, \bar d_i + \lambda_k\, \overline{\mathrm{sh}}_i + \mu\, \mathrm{pd}_{d(i)} + \tfrac12\big(\varepsilon_{i,\mathrm{B}} + \varepsilon_{i,\mathrm{A}}\big), \label{eq:between}
\end{align}
with $\Delta d_i$ and $\Delta \mathrm{sh}_i$ the changes in premeasured style mismatch and competence shortfall from A to B, $\bar d_i$ and $\overline{\mathrm{sh}}_i$ their means, $g_i$ the group indicator, $\tilde a_i$ the baseline's zero-mean part, $\alpha$ an intercept absorbing $\tfrac12 \gamma$, $u_d$ the dyad effect the simulation adds to \eqref{eq:family} and $\mathrm{pd}_d$ the partners' style distance. The group indicator carries the baseline gap $\Delta a$; without it, where the rule sets are not mirror images of each other, the group difference in mean mismatch would be absorbed by $\mu$ and the shared $\lambda$'s, and the real study's rule sets will not be mirror images.

Two assumptions make the decomposition a likelihood, and balance supplies neither. The residual variances are equal across conditions and the errors jointly normal: then
\[ \operatorname{Cov}\!\Big(\varepsilon_{i,\mathrm{B}} - \varepsilon_{i,\mathrm{A}},\ \tfrac12\big(\varepsilon_{i,\mathrm{B}} + \varepsilon_{i,\mathrm{A}}\big)\Big) \;=\; \tfrac12\big(\operatorname{Var}\varepsilon_{i,\mathrm{B}} - \operatorname{Var}\varepsilon_{i,\mathrm{A}}\big) \;=\; 0, \]
and under normality the difference and the mean are independent, so that \eqref{eq:within} identifies $\gamma$, $\lambda_s$ and $\lambda_k$ with the baseline, the dyad effect and the partner term cancelled, and \eqref{eq:between} identifies $\mu$ and $\Delta a$; the robustness study reports what unequal variances do. Second, the dyad effect cancels from the differences but not from the means, whose two partners it correlates with covariance $\sigma_u^2$ ($0.16$ at the simulation's setting), so the $N$ means are not independent observations. The $N$ differences are independent of each other and of the means; within each dyad the two partners' means, with compound-symmetric covariance, are rotated into their sum and their difference, each divided by $\sqrt{2}$, which are uncorrelated and under normality independent, with variances $v + \sigma_u^2$ and $v - \sigma_u^2$ for $v$ the variance of a person mean about its expectation, and the design of each rotated part is the between design rotated the same way. The likelihood is therefore a product over $2N$ independent Gaussian terms, the $N$ differences with variance $2\sigma_\varepsilon^2$, the $N/2$ rotated sums and the $N/2$ rotated differences, each part with its own variance and the coefficients shared between the parts; the maximum-likelihood fit is iterated weighted least squares, run for forty iterations, well past convergence. Models are compared by $-2\log L + k \log(2N)$, with $k$ the number of coefficients plus the three variances and $2N$ the number of independent terms, which is also the number of effort observations. The identification of $\mu$ rests on the baselines $\tilde a_i$ being mean-independent of $\mathrm{pd}_{d(i)}$ given $g_i$ and the measured profiles, which the simulation satisfies by construction except in the robustness scenario built to violate it.

The parameter-recovery study fits each equation by least squares with cluster-robust standard errors by dyad and tests $(\lambda_s, \lambda_k)$ jointly by a Wald statistic on the clustered covariance. The model-recovery study compares the four restrictions by the criterion above and by the predictive log score of the held-out dyads under the fitted Gaussian parts, six folds. The withheld-rule-set study fits \eqref{eq:within} on the training dyads, predicts $y_{i,\mathrm{C}} - y_{i,\mathrm{A}}$ for the held-out dyads from the fitted $\lambda$'s and the premeasured mismatch under C, centres the prediction and the observation within the held-out set, and compares the root-mean-square error with the noise floor $\sqrt{2}\,\sigma_\varepsilon\sqrt{1 - 1/n_f}$, $n_f$ the number of held-out participants.

\section{Preregistration items}
\label{app:prereg}

The preregistration fixes, before any participant is run: the instruments for $s_i$ and $k_i$ and their scales; the two rule sets' texts, their coded $s^{*}_n$, $\theta_n$ and $w_n$, and the rater protocol; the distance $d$; the effort measures, subjective and behavioural, and the expected and perceived empathy reports; the practice-to-criterion rule and the adherence coding; the counterbalancing; the primary analysis of \eqref{eq:within} with dyad-clustered errors, its primary test the joint Wald test of $(\lambda_s, \lambda_k)$ at $5\%$, the secondary one-sided tests of each coefficient in the predicted direction, and the smallest effects of interest for the equivalence tests; the period and carryover terms; the covariates of training time, practice errors and practice effort; the actor--partner model for levels, with its identifying assumption stated; the model comparison by the Bayesian information criterion on the exact likelihood; the mediation analysis with its sensitivity analysis; the three failure conditions of \autoref{sec:design}; and the sample size, $90$ mixed dyads if the full family is to be compared and $60$ if the primary prediction alone is at issue, at the simulated effect sizes, which are to be replaced by pilot estimates where available. Before the preregistration, a pilot establishes whether the profiles and the norm requirements can be scored on shared scales, whether participants can follow both rule sets to criterion and at what cost in training, whether the effort measures track the intended construct against performance and adherence, and the reliabilities of the empathy reports and of the rule-set codings.


\section{The source passages}
\label{app:source}

The passages below are quoted as written in the larger work \cite{Aethus}, from the section ``The Aethic Mindset and Autistic Savants'' and its subsection ``Illustrating the Bias of Perspective in Diagnosing Logic'', and from the paragraph ``The decisive fingerprint: it toggles'' of the section ``Aethic Reasoning Against the Modal Structure of Human Language''. The source dates the two principles to a short video the author recorded on June 8, 2024, at 2:39 PM ET.

\paragraph{The Branch Principle, and its motivation.}
\begin{quote}
``Ontologically neurodivergent people are better understood as a collection of categorical branches and `allistic' is a single branch, such that two random branches are typically as different from one another as they each are to a third.''
\end{quote}
\begin{quote}
``A better model for understanding autistic minds than the autistic-allistic binary is more of a tree model, such that allistic people are one branch alone and the rest have been all categorized as `autistic'. Inadvertently speaking, then, this proposes that an autism diagnosis is not actually diagnosing a person's traits to the level of specificity of an allistic person (as would assign them a particular branch) but instead broadly classifies them as `non-allistic' for their lacking allistic traits. In reality, though, each branch comes with a different set of cognitive intuitions, which are separate enough so that two autistic people drawn from two different branches tend to be as different from one another at the cognitive level as either of them are from allistic people.''
\end{quote}

\paragraph{The Principle of Symmetry, and its intended content.}
\begin{quote}
``You can always redraw the conventional category boundaries in psychology so that allistic and a given neurodivergent branch `trade places phenomenologically', in the sense that the exact behavioral diagnosis, as an Aethus, is made to reverse isomorphically.''
\end{quote}
``As an Aethus'' is glossed at the source as ``a partial information state about that behavioral diagnosis rather than an exhaustive description''. The principle's intended content is given in the same passage with the author's own case as the instance, offered in this paper as motivation only:
\begin{quote}
``What's fundamentally wrong about this is that it fundamentally misunderstands what's actually going on: what's going on is utter, complete symmetry between cognitive branches. That is, for allistic people versus a given branch of autistic people, there will be roughly half-and-half instances across traits for which the autistic person is infantile versus for which the allistic person is infantile. Using myself as an example, at age twenty I still had the emotional maturity of a thirteen year old, but at age thirteen I already had the mathematical maturity of at least a twenty year old. Had my branch been the reference instead, allistic people would look just as profoundly infantile in so many different ways to how I looked in social, emotional, and life skills.''
\end{quote}

\paragraph{The random-partition illustration.}
\begin{quote}
``This is also where categorizing all the non-allistic branches as one `autism' is measurably problematic: you can find an autistic branch for which any trait is infantile, thereby seeming to confirm your prior that `look, autistic people are infantile in this way'. Running the same logic on a random partition where we put half the branches, including the allistic one, in one category while the rest are placed in the other, and we could just as easily misinterpret allistic people as infantile in any conceivable way by finding it in another branch of the categorization and then mistakenly importing it to an assessment of allistic people.''
\end{quote}

\paragraph{The seat of the cost.} From a footnote of the same section:
\begin{quote}
``[N]otice how many autistic people themselves are described as taking a script to allistic conversations. Once again, notice how if the shoe were on the other foot, for any given non-allistic cognition branch allistic people practicing its topics would be interpretable as using a script.''
\end{quote}

\paragraph{The toggle.}
\begin{quote}
``One further fingerprint is of a different kind, and it is decisive against the strongest rival reading: that the convention is not a convention at all but a fixed cognitive limitation. A limitation does not toggle. The switching just described---an adult moving, on demand, between the attached picture and the conventional one within a single day, by context---is a \emph{dissociability} result: the convention is a mode the platform runs, selectable, not a ceiling the platform imposes.''
\end{quote}

\end{document}
